\documentclass{article}
\usepackage[T1]{fontenc}
\usepackage{lmodern}
\usepackage{graphicx}
\usepackage{amsmath,amssymb}
\usepackage[a4paper,margin=2cm]{geometry}
\usepackage[absolute,overlay]{textpos}
\setlength{\TPHorizModule}{1mm}
\setlength{\TPVertModule}{1mm}
\usepackage[indonesian]{babel}
\usepackage{hyperref}
\setlength{\emergencystretch}{2em}
% unit: Halaman 65

\begin{document}

% === Halaman 65 (llm) ===

\section*{Jaringan Feistel}

\begin{itemize}
    \item Beberapa cipher blok seperti DES, GOST, Feal, Lucifer, RC5, RC6, dan lain-lain menggunakan struktur \textit{enciphering} pada setiap putaran yang dinamakan \textbf{jaringan Feistel} (\textit{Feistel Network}).
    \item Dalam hal ini, blok plainteks dibagi menjadi dua bagian, dan pada masing-masing bagian dilakukan proses transformasi menjadi upa-bagian pada putaran selanjutnya.
    \item Struktur ini bersifat \textit{reversible}, yaitu untuk melakukan dekripsi maka jaringan Feistel dieksekusi dari "bawah" ke "atas".
\end{itemize}

\end{document}
